% GENERATED FILE. Edit canonical lessons, then run npm run generate:books.
% canonical-source-hash: fnv1a64:8bdb3cd2b6f60a62
% canonical-lessons: RU-C146-bronirovat, RU-C146-povorachivat, RU-C146-peresekat, RU-C146-ostanavlivatsya, RU-C146-parkovat

\chapter{To book, To turn, To cross, To stop, To park}
\label{ch:ru-to-book-to-turn-to-cross-to-stop-to-park}
\hlchaptermodality{146}

\begin{chapteropening}
\textbf{By the end of this chapter:} \emph{I can say to book, to turn, to cross, to stop and to park.}

Complete the last lesson of chapter 146: I can say to book, to turn, to cross, to stop and to park.
\end{chapteropening}

\section[\texorpdfstring{\ru{бронировать} (bronírovat)}{бронировать (bronírovat)}]{\ru{бронировать} (bronírovat) — to book, to reserve}

\label{lesson:RU-C146-bronirovat}



\begin{quote}
Before the new one: say the Russian for to download, then the Russian for to be beaten.
\end{quote}

\subsection*{You'll want to know: \ru{бронировать}}
\textbf{\ru{бронировать}} (\emph{bronírovat}) — \textquotedblleft{}to book, to reserve\textquotedblright{}.

\begin{grammarlens}[title={What you've built}]
One more word for this chapter.
\end{grammarlens}

\subsection*{Your turn}
\begin{itemize}
\raggedright
  \item \emph{Say it:} \emph{bronírovat}
  \item \emph{Say it:} \emph{bronírovat}, once more
  \item \emph{Say it:} say \emph{bronírovat}
  \item \emph{Recall:} say the Russian for to download, then the Russian for to be beaten, then say \emph{bronírovat} again
\end{itemize}

\begin{tcolorbox}[breakable,colback=teal!4,colframe=teal!35!black,title={Before you move on}]
What does \ru{бронировать} mean? (\textquotedblleft{}To book, to reserve\textquotedblright{}.) Say it once more.
\end{tcolorbox}
\section[\texorpdfstring{\ru{поворачивать} (povoráchivat)}{поворачивать (povoráchivat)}]{\ru{поворачивать} (povoráchivat) — to turn}

\label{lesson:RU-C146-povorachivat}



\begin{quote}
Before the new one: say the Russian for to be beaten, then the Russian for to book.
\end{quote}

\subsection*{You'll want to know: \ru{поворачивать}}
\textbf{\ru{поворачивать}} (\emph{povoráchivat}) — \textquotedblleft{}to turn\textquotedblright{}.

\begin{grammarlens}[title={What you've built}]
One more word for this chapter.
\end{grammarlens}

\subsection*{Your turn}
\begin{itemize}
\raggedright
  \item \emph{Say it:} \emph{povoráchivat}
  \item \emph{Say it:} \emph{povoráchivat}, once more
  \item \emph{Say it:} say \emph{povoráchivat}
  \item \emph{Recall:} say the Russian for to be beaten, then the Russian for to book, then say \emph{povoráchivat} again
\end{itemize}

\begin{tcolorbox}[breakable,colback=teal!4,colframe=teal!35!black,title={Before you move on}]
What does \ru{поворачивать} mean? (\textquotedblleft{}To turn\textquotedblright{}.) Say it once more.
\end{tcolorbox}
\section[\texorpdfstring{\ru{пересекать} (peresekát)}{пересекать (peresekát)}]{\ru{пересекать} (peresekát) — to cross}

\label{lesson:RU-C146-peresekat}



\begin{quote}
Before the new one: say the Russian for to book, then the Russian for to turn.
\end{quote}

\subsection*{You'll want to know: \ru{пересекать}}
\textbf{\ru{пересекать}} (\emph{peresekát}) — \textquotedblleft{}to cross\textquotedblright{}.

\begin{grammarlens}[title={What you've built}]
One more word for this chapter.
\end{grammarlens}

\subsection*{Your turn}
\begin{itemize}
\raggedright
  \item \emph{Say it:} \emph{peresekát}
  \item \emph{Say it:} \emph{peresekát}, once more
  \item \emph{Say it:} say \emph{peresekát}
  \item \emph{Recall:} say the Russian for to book, then the Russian for to turn, then say \emph{peresekát} again
\end{itemize}

\begin{tcolorbox}[breakable,colback=teal!4,colframe=teal!35!black,title={Before you move on}]
What does \ru{пересекать} mean? (\textquotedblleft{}To cross\textquotedblright{}.) Say it once more.
\end{tcolorbox}
\section[\texorpdfstring{\ru{останавливаться} (ostanávlivatsya)}{останавливаться (ostanávlivatsya)}]{\ru{останавливаться} (ostanávlivatsya) — to stop}

\label{lesson:RU-C146-ostanavlivatsya}



\begin{quote}
Before the new one: say the Russian for to turn, then the Russian for to cross.
\end{quote}

\subsection*{You'll want to know: \ru{останавливаться}}
\textbf{\ru{останавливаться}} (\emph{ostanávlivatsya}) — \textquotedblleft{}to stop\textquotedblright{}.

\begin{grammarlens}[title={What you've built}]
One more word for this chapter.
\end{grammarlens}

\subsection*{Your turn}
\begin{itemize}
\raggedright
  \item \emph{Say it:} \emph{ostanávlivatsya}
  \item \emph{Say it:} \emph{ostanávlivatsya}, once more
  \item \emph{Say it:} say \emph{ostanávlivatsya}
  \item \emph{Recall:} say the Russian for to turn, then the Russian for to cross, then say \emph{ostanávlivatsya} again
\end{itemize}

\begin{tcolorbox}[breakable,colback=teal!4,colframe=teal!35!black,title={Before you move on}]
What does \ru{останавливаться} mean? (\textquotedblleft{}To stop\textquotedblright{}.) Say it once more.
\end{tcolorbox}
\section[\texorpdfstring{\ru{парковать} (parkovát)}{парковать (parkovát)}]{\ru{парковать} (parkovát) — to park}

\label{lesson:RU-C146-parkovat}



\begin{quote}
Before the new one: say the Russian for to cross, then the Russian for to stop.
\end{quote}

\subsection*{You'll want to know: \ru{парковать}}
\textbf{\ru{парковать}} (\emph{parkovát}) — \textquotedblleft{}to park\textquotedblright{}.

\begin{grammarlens}[title={What you've built}]
One more word for this chapter.
\end{grammarlens}

\subsection*{Your turn}
\begin{itemize}
\raggedright
  \item \emph{Say it:} \emph{parkovát}
  \item \emph{Say it:} \emph{parkovát}, once more
  \item \emph{Say it:} say \emph{parkovát}
  \item \emph{Recall:} say the Russian for to cross, then the Russian for to stop, then say \emph{parkovát} again
\end{itemize}

\begin{tcolorbox}[breakable,colback=teal!4,colframe=teal!35!black,title={Before you move on}]
What does \ru{парковать} mean? (\textquotedblleft{}To park\textquotedblright{}.) Say it once more.
\end{tcolorbox}
